\documentclass[11pt]{article}
\usepackage[margin=1in]{geometry}
\usepackage{amsmath, amssymb, bm}
\usepackage{graphicx}
\usepackage[T1]{fontenc}
\usepackage[utf8]{inputenc}
\usepackage{lmodern}
\usepackage{microtype}
\usepackage{enumitem}
\setlist[itemize]{leftmargin=1.5em}

\title{SAVLG-CBM Ablation Methodology}
\author{}
\date{\today}

\begin{document}
\maketitle

\section{Pipeline}
For an input image $x$, the SAVLG model produces a global concept branch and a
spatial concept branch.

\subsection{Backbone Features}
Let the frozen visual backbone produce:
\[
\mathbf{f}_4 \in \mathbb{R}^{C_4 \times H_4 \times W_4},
\qquad
\mathbf{f}_5 \in \mathbb{R}^{C_5 \times H_5 \times W_5}.
\]
For the ResNet-50 MMPretrain CUB setup with $448 \times 448$ inputs:
\[
H_4 \times W_4 = 28 \times 28,
\qquad
H_5 \times W_5 = 14 \times 14.
\]

\subsection{Global Branch}
The global branch uses the conv5 feature map. Global average pooling gives
\[
\bar{\mathbf{f}}_5 = \operatorname{GAP}(\mathbf{f}_5),
\]
and the global concept head produces
\[
\mathbf{c}_{\mathrm{global}} = g(\bar{\mathbf{f}}_5) \in \mathbb{R}^{K},
\]
where $K$ is the number of concepts.

\subsection{Spatial Branch}
The spatial branch predicts concept maps
\[
\mathbf{M} \in \mathbb{R}^{K \times H \times W}.
\]
The exact branch depends on the ablation:

\paragraph{conv5 branch.}
\[
\mathbf{M} = s(\mathbf{f}_5),
\qquad
H \times W = 14 \times 14.
\]

\paragraph{conv4 branch.}
\[
\mathbf{M} = s(\mathbf{f}_4),
\qquad
H \times W = 28 \times 28.
\]

\paragraph{conv4+conv5 multiscale branch.}
The two stages are fused first:
\[
\tilde{\mathbf{f}}_5 = \operatorname{Up}\!\left(W_5 * \mathbf{f}_5\right),
\qquad
\tilde{\mathbf{f}}_4 = W_4 * \mathbf{f}_4,
\]
\[
\mathbf{f}_{\mathrm{fused}} = \operatorname{ReLU}\!\left(\tilde{\mathbf{f}}_4 + \tilde{\mathbf{f}}_5\right),
\]
then
\[
\mathbf{M} = s(\mathbf{f}_{\mathrm{fused}}),
\qquad
H \times W = 28 \times 28.
\]
So the multiscale branch does \emph{early fusion of features}, not separate
losses on separate conv4 and conv5 maps.

\subsection{Spatial Logit Pooling}
The spatial maps are pooled into spatial concept logits:
\[
\mathbf{c}_{\mathrm{spatial}} = p(\mathbf{M}) \in \mathbb{R}^{K},
\]
where the residual coupling path uses LSE pooling in the current ablations.

\subsection{Residual Concept Fusion}
The final concept vector seen by the classifier is
\[
\mathbf{c}_{\mathrm{final}}
=
\mathbf{c}_{\mathrm{global}}
+
\alpha \, \mathbf{c}_{\mathrm{spatial}},
\]
where $\alpha$ is the residual spatial coupling coefficient.

\paragraph{Important detail.}
In the current implementation, the image-level concept loss is applied to
$\mathbf{c}_{\mathrm{final}}$, not just $\mathbf{c}_{\mathrm{global}}$.
So the spatial branch affects training through the fused concept logits.

\section{Loss Used in the 5-Run Ablation Study}
The ablation study only used the following core terms:
\[
\mathcal{L}
=
\lambda_{\mathrm{global}} \mathcal{L}_{\mathrm{global}}
+
\lambda_{\mathrm{align}} \mathcal{L}_{\mathrm{align}}
+
\lambda_{\mathrm{outside}} \mathcal{L}_{\mathrm{outside}}
+
\lambda_{\mathrm{coverage}} \mathcal{L}_{\mathrm{coverage}}
+
\lambda_{\mathrm{dice}} \mathcal{L}_{\mathrm{dice}}.
\]

In most runs:
\[
\lambda_{\mathrm{dice}} = 0.
\]
Only the Dice ablation enabled the final term.

\section{Loss Terms}
\subsection{1. Global Concept BCE on Fused Logits}
Let $\mathbf{y} \in \{0,1\}^{K}$ denote the image-level concept targets. The
global concept supervision used in the ablations is plain BCE on the
\emph{fused} logits:
\[
\mathcal{L}_{\mathrm{global}}
=
\frac{1}{K}\sum_{k=1}^{K}
\operatorname{BCEWithLogits}\!\left(c_{\mathrm{final},k}, y_k\right).
\]

\paragraph{Interpretation.}
This term trains the final concept vector used by the sparse classifier. Because
$\mathbf{c}_{\mathrm{final}}$ contains the residual spatial contribution, the
spatial branch participates in image-level concept learning.

\subsection{2. Soft-Align KL Loss on Spatial Maps}
For a supervised concept $k$, let:
\[
\mathbf{m}_k \in \mathbb{R}^{H' \times W'}
\]
be the interpolated predicted concept map logits and
\[
\mathbf{t}_k \in [0,1]^{H' \times W'}
\]
be the target soft box / mask target.

The target is normalized into a spatial distribution:
\[
\mathbf{q}_k
=
\frac{\mathbf{t}_k}{\sum_u t_{k,u}}.
\]
The prediction becomes a log-distribution:
\[
\log \mathbf{p}_k
=
\log \operatorname{softmax}(\mathbf{m}_k).
\]
Then the local alignment loss is
\[
\mathcal{L}_{\mathrm{align}}
=
D_{\mathrm{KL}}(\mathbf{q}_k \,\|\, \mathbf{p}_k).
\]

\paragraph{Interpretation.}
This encourages the spatial map to put probability mass where the target region
lies, without forcing hard binary pixel classification.

\subsection{Example: Bounding Box to $t_k$ and $q_k$}
Figure~\ref{fig:target-mask-rasterization} shows the geometric picture: a
concept-level bounding box is rasterized onto the spatial supervision grid,
producing the soft target mask $t_k$.

\begin{figure}[h]
    \centering
    \includegraphics[width=0.92\linewidth]{docs/figures/savlg_target_mask_example.pdf}
    \caption{A concept bounding box in the image plane is rasterized onto the
    spatial supervision grid. In \texttt{soft\_box} mode, each entry of
    $t_k$ is the fractional occupancy of the corresponding grid cell.}
    \label{fig:target-mask-rasterization}
\end{figure}

Suppose an image-level annotation provides a concept box that covers the middle
of the object, and suppose the supervision grid is $4 \times 4$. In the
\texttt{soft\_box} target mode, each entry of $t_k$ is the fractional occupancy
of the corresponding grid cell by the bounding box. A simple example is:
\[
\mathbf{t}_k
=
\begin{bmatrix}
0 & 0 & 0 & 0 \\
0 & 0.25 & 0.50 & 0 \\
0 & 0.50 & 1.00 & 0 \\
0 & 0 & 0 & 0
\end{bmatrix}.
\]
The total target mass is
\[
\sum_u t_{k,u} = 2.25.
\]
Therefore the normalized soft target distribution used by the KL term is
\[
\mathbf{q}_k
=
\frac{\mathbf{t}_k}{2.25}
=
\begin{bmatrix}
0 & 0 & 0 & 0 \\
0 & 0.1111 & 0.2222 & 0 \\
0 & 0.2222 & 0.4444 & 0 \\
0 & 0 & 0 & 0
\end{bmatrix}.
\]

\paragraph{Interpretation.}
Here, $t_k$ is still a spatial support target in the original mask units,
whereas $q_k$ is the probability distribution used by the soft-align KL loss.
Large cells inside the box receive more target probability mass than partially
covered cells near the boundary. In the implementation, small clamps are used
for numerical safety, but conceptually this normalization is just
$q_k = t_k / \sum_u t_{k,u}$.

\subsection{Why These Spatial Terms Are Needed}
The three main spatial terms used in the ablation study play different roles.
They are easiest to understand by considering three failure modes of a concept
map:

\paragraph{Failure mode A: correct region, but mass placed too unevenly.}
The model may activate inside the target box, but concentrate almost all mass in
one corner or one very small hotspot. In that case, the map is ``inside'' the
box, but it does not match the intended spatial support well.

\paragraph{Failure mode B: activation spills outside the target.}
The model may put substantial probability mass on background or neighboring
parts. This can still produce a nontrivial overlap score, but the map is not
cleanly grounded.

\paragraph{Failure mode C: tiny peak inside a large target.}
The model may place a sharp peak inside the box and satisfy a weak containment
style notion, yet fail to cover the full target extent. This is especially bad
for IoU-oriented evaluation.

These failure modes motivate the three terms below:

\begin{itemize}
    \item \textbf{Soft-align KL:} aligns the predicted spatial probability map to
    the normalized target distribution $q_k$, so the map does not merely land
    somewhere inside the box but matches the target support more faithfully.
    \item \textbf{Outside-mass penalty:} explicitly suppresses probability mass
    outside the target support, reducing spillover onto irrelevant regions.
    \item \textbf{Coverage penalty:} rewards activating across the target region
    rather than producing a tiny discriminative hotspot.
\end{itemize}

\paragraph{Combined effect.}
Soft-align shapes \emph{where} mass should go, outside penalty removes
\emph{where mass should not go}, and coverage encourages \emph{how much of the
target region should be covered}. In practice, this combination was more useful
than using any single term alone.

\subsection{3. Outside-Mass Penalty}
Let
\[
\tilde{\mathbf{p}}_k = \operatorname{softmax}(\mathbf{m}_k)
\]
over spatial locations. The outside-mass penalty is
\[
\mathcal{L}_{\mathrm{outside}}
=
\sum_u \tilde{p}_{k,u}(1 - t_{k,u}).
\]

\paragraph{Interpretation.}
This penalizes probability mass outside the target support and reduces spillover.

\subsection{4. Coverage Penalty}
Let
\[
\hat{\mathbf{p}}_k = \sigma(\mathbf{m}_k).
\]
For one supervised concept $k$, coverage over the target region is
\[
\operatorname{coverage}_k
=
\frac{
\sum_u \hat{p}_{k,u} t_{k,u}
}{
\sum_u t_{k,u} + \varepsilon
}.
\]
The corresponding per-concept penalty is
\[
\mathcal{L}_{\mathrm{coverage}}^{(k)}
=
1 - \operatorname{coverage}_k.
\]

The full training loss averages this quantity over the supervised concepts in an
image, and then averages again over the batch:
\[
\mathcal{L}_{\mathrm{coverage}}
=
\frac{1}{B}\sum_{b=1}^{B}
\left(
\frac{
\sum_{k \in \mathcal{K}_b}
w_{b,k}\,
\mathcal{L}_{\mathrm{coverage}}^{(b,k)}
}{
\sum_{k \in \mathcal{K}_b} w_{b,k}
}
\right),
\]
where $\mathcal{K}_b$ is the set of supervised concepts for image $b$.

\paragraph{Interpretation.}
This pushes the map to cover more of the target region rather than only peaking
inside it. Because the numerator is
\[
\sum_u \hat{p}_{k,u} t_{k,u},
\]
locations with larger target values $t_{k,u}$ contribute more to the covered
mass. So yes: if a location has high $t_{k,u}$, the loss encourages high
$\hat{p}_{k,u}$ there. In the simplest binary case, this term rewards having
high sigmoid probability across the whole target support rather than only at one
small hotspot.

\subsection{5. Dice Loss}
Only the Dice ablation used this term. With sigmoid probabilities:
\[
\hat{\mathbf{p}}_k = \sigma(\mathbf{m}_k),
\]
the Dice loss is
\[
\mathcal{L}_{\mathrm{dice}}
=
1 -
\frac{
2 \sum_u \hat{p}_{k,u} t_{k,u} + \varepsilon
}{
\sum_u \hat{p}_{k,u} + \sum_u t_{k,u} + \varepsilon
}.
\]

\paragraph{Interpretation.}
Dice directly rewards overlap between predicted support and target support.

\section{Ablation-Specific Loss Stacks}
\paragraph{conv5 + cov025.}
\[
\mathcal{L}
=
\mathcal{L}_{\mathrm{global}}
+
0.25\,\mathcal{L}_{\mathrm{align}}
+
0.25\,\mathcal{L}_{\mathrm{outside}}
+
0.25\,\mathcal{L}_{\mathrm{coverage}}.
\]

\paragraph{conv5 + cov050.}
Same as above, but with
\[
\lambda_{\mathrm{coverage}} = 0.5.
\]

\paragraph{conv5 + dice010.}
Same as `conv5 + cov025`, plus
\[
\lambda_{\mathrm{dice}} = 0.1.
\]

\paragraph{conv4.}
Same loss stack as `conv5 + cov025`, but with the spatial branch operating on
conv4 features.

\paragraph{conv4+conv5.}
Same loss stack as `conv5 + cov025`, but with the spatial branch operating on
the fused conv4+conv5 feature map.

\section{Summary}
The 5-run ablation study varied two things:
\begin{itemize}
    \item the spatial feature source: conv5, conv4, or conv4+conv5,
    \item the local overlap regularization: baseline coverage, higher coverage, or Dice.
\end{itemize}

The key design point is:
\begin{itemize}
    \item concept supervision is applied to the fused logits $\mathbf{c}_{\mathrm{final}}$,
    \item localization supervision is applied directly to the spatial maps $\mathbf{M}$.
\end{itemize}

\end{document}
